\section{One-Shot Results and Matched-Effort Comparison}
\label{app:oneshot}

\subsection{One-Shot Results Across Reasoning Efforts}

\begin{table}[!t]
\centering\footnotesize
\caption{One-shot results for every model and supported reasoning effort, Layout and shading follow Table~\ref{tab:main}; shaded rows mark each model's highest-scoring one-shot effort. The one-shot setting receives no performance feedback.}
\label{tab:oneshot}
\setlength{\tabcolsep}{3.4pt}
\begin{tabular}{@{}llrrrrrrrr@{}}
\toprule
Model & Effort & $S$ & $S_{\mathrm{in}}$ & $S_{\mathrm{out}}$ & ARR (\%) & Excess (pp) & RankIC & Sharpe & MDD (\%) \\
\midrule
\modellogo{claude}\ModelClaudeOpus & low & 35.04 & 43.21 & 31.53 & $-4.5$ & $-21.6$ & $-0.011$ & 0.00 & $-33.0$ \\
\rowcolor{gray!14} & medium & \textbf{42.61} & 47.19 & \textbf{40.65} & 5.6 & $-11.5$ & $-0.001$ & 0.34 & $\mathbf{-21.0}$ \\
 & high & 36.07 & 39.68 & 34.52 & $-7.5$ & $-24.5$ & 0.020 & $-0.14$ & $-29.7$ \\
 & xhigh & 39.13 & 49.05 & 34.88 & $-9.9$ & $-26.9$ & 0.040 & $-0.20$ & $-35.0$ \\
 & max & 39.70 & 42.07 & 38.69 & $-3.7$ & $-20.8$ & \textbf{0.046} & 0.03 & $-35.9$ \\
\midrule
\rowcolor{gray!14}\modellogo{openai}\ModelTerra & none & 42.14 & 46.56 & 40.25 & \textbf{9.5} & $\mathbf{-7.6}$ & $-0.010$ & \textbf{0.44} & $-27.8$ \\
 & low & 27.82 & 26.59 & 28.35 & $-10.3$ & $-27.4$ & $-0.007$ & $-0.24$ & $-25.7$ \\
 & medium & 21.01 & 20.44 & 21.25 & $-32.4$ & $-49.4$ & $-0.010$ & $-1.51$ & $-34.1$ \\
 & high & 40.61 & 42.88 & 39.64 & 6.2 & $-10.9$ & $-0.001$ & 0.35 & $-27.5$ \\
 & xhigh & 34.21 & 42.94 & 30.46 & $-9.8$ & $-26.8$ & 0.000 & $-0.17$ & $-33.6$ \\
 & max & 20.98 & 20.98 & 20.97 & $-32.3$ & $-49.3$ & $-0.009$ & $-1.60$ & $-32.6$ \\
\midrule
\modellogo{claude}\ModelClaudeSonnet & low & 37.75 & 42.41 & 35.76 & 1.1 & $-16.0$ & $-0.007$ & 0.19 & $-30.3$ \\
 & medium & 33.82 & 41.81 & 30.39 & $-11.5$ & $-28.6$ & 0.002 & $-0.21$ & $-35.5$ \\
\rowcolor{gray!14} & high & 40.24 & 44.28 & 38.50 & 1.4 & $-15.6$ & 0.010 & 0.20 & $-32.1$ \\
 & xhigh & 37.05 & 47.90 & 32.40 & $-14.3$ & $-31.3$ & 0.029 & $-0.33$ & $-36.0$ \\
 & max & 33.89 & 41.19 & 30.77 & $-14.8$ & $-31.9$ & 0.018 & $-0.33$ & $-36.9$ \\
\midrule
\rowcolor{gray!14}\modellogo{glm}\ModelGLM & high & 38.02 & \textbf{52.54} & 31.79 & $-4.2$ & $-21.3$ & $-0.023$ & 0.05 & $-37.3$ \\
 & max & 31.46 & 34.70 & 30.07 & $-15.4$ & $-32.4$ & 0.022 & $-0.51$ & $-32.1$ \\
\midrule
\modellogo{gemini}\ModelGemini & low & 27.75 & 21.26 & 30.53 & $-12.1$ & $-29.1$ & 0.005 & $-0.41$ & $-23.5$ \\
\rowcolor{gray!14} & medium & 35.91 & 35.62 & 36.04 & 1.6 & $-15.4$ & $-0.006$ & 0.19 & $-23.8$ \\
 & high & 30.26 & 25.93 & 32.11 & $-5.4$ & $-22.5$ & $-0.003$ & $-0.06$ & $-24.6$ \\
\midrule
\modellogo{openai}\ModelAstra & low & 30.08 & 35.05 & 27.95 & $-11.2$ & $-28.3$ & $-0.010$ & $-0.28$ & $-30.3$ \\
 & medium & 31.02 & 36.73 & 28.57 & $-9.7$ & $-26.7$ & $-0.012$ & $-0.20$ & $-30.3$ \\
\rowcolor{gray!14} & high & 35.22 & 38.27 & 33.92 & $-4.8$ & $-21.9$ & 0.007 & $-0.03$ & $-30.9$ \\
 & xhigh & 30.30 & 38.46 & 26.80 & $-16.7$ & $-33.8$ & $-0.002$ & $-0.48$ & $-34.7$ \\
 & max & 31.47 & 40.13 & 27.75 & $-18.3$ & $-35.3$ & 0.007 & $-0.54$ & $-34.8$ \\
\midrule
\rowcolor{gray!14}\modellogo{openai}\ModelSol & low & 32.15 & 38.29 & 29.52 & $-7.6$ & $-24.7$ & $-0.014$ & $-0.08$ & $-32.2$ \\
 & medium & 22.60 & 22.62 & 22.59 & $-31.4$ & $-48.4$ & $-0.006$ & $-1.31$ & $-34.1$ \\
 & high & 23.15 & 21.07 & 24.04 & $-17.4$ & $-34.4$ & $-0.020$ & $-0.51$ & $-29.6$ \\
 & xhigh & 23.29 & 23.84 & 23.05 & $-19.0$ & $-36.1$ & $-0.022$ & $-0.60$ & $-29.6$ \\
 & max & 22.31 & 24.08 & 21.55 & $-23.7$ & $-40.7$ & $-0.018$ & $-0.98$ & $-32.5$ \\
\midrule
\modellogo{kimi}\ModelKimi & low & 31.58 & 38.58 & 28.58 & $-13.4$ & $-30.5$ & 0.000 & $-0.34$ & $-33.8$ \\
\rowcolor{gray!14} & high & 32.01 & 37.89 & 29.48 & $-20.5$ & $-37.6$ & 0.038 & $-0.84$ & $-38.6$ \\
 & max & 31.65 & 40.87 & 27.70 & $-20.3$ & $-37.3$ & 0.016 & $-0.86$ & $-35.1$ \\
\midrule
\modellogo{deepseek}\ModelDeepSeek & none & 20.48 & 22.01 & 19.82 & $-34.1$ & $-51.1$ & $-0.015$ & $-1.79$ & $-37.5$ \\
 & low & 29.36 & 34.21 & 27.28 & $-17.4$ & $-34.4$ & 0.001 & $-0.49$ & $-33.3$ \\
 & high & 29.63 & 37.51 & 26.25 & $-25.4$ & $-42.4$ & 0.010 & $-0.92$ & $-39.5$ \\
\rowcolor{gray!14} & max & 29.92 & 35.95 & 27.34 & $-11.9$ & $-29.0$ & $-0.008$ & $-0.42$ & $-26.7$ \\
\bottomrule
\end{tabular}

\end{table}

Table~\ref{tab:oneshot} reports the 38 one-shot runs in the layout of Table~\ref{tab:main}. The median configuration scores 31.61, and 4 of the 38 exceed 40 points, against 21 in the iterative setting. Six configurations have positive sealed-period ARR, and none exceeds the 17.1\% market.

\noindent\textbf{Without feedback, reasoning effort is just as inconsistent as with it.} For eight of the nine models the best effort is not the highest one supported (Table~\ref{tab:oneshot}). Terra scores 42.14 with no reasoning and 20.98 at max, a 21.16-point range that exceeds every iterative range, and Sol is best at low (32.15) and lowest at max (22.31). Only DeepSeek rises with effort, from 20.48 with no reasoning to 29.92 at max. With no feedback loop to correct a poor draft, each cell reflects a single draft, so the scores swing more with effort.

\subsection{Matched-Effort Comparison}

\noindent\textbf{Iteration helps at matched efforts as well.} At each model's highest supported effort (Gemini: high; others: max), the iterative run outscores the one-shot run for all nine models, with mean scores of 40.88 against 30.18, and has the higher ARR in six of the nine pairs. Across all 38 matched model--effort cells, the iterative run has the higher score in 34, the higher ARR in 26, and the higher RankIC in 28. Matching efforts gives a larger gap than pairing best configurations (Table~\ref{tab:paired}) because one-shot scores fall with effort more often, so a one-shot run at max effort is rarely its model's best.
